% report.tex
%
% Purpose:
%   Give a self-contained theoretical investigation of Conjecture 1 in the
%   TEP draft: a single-cell Helmholtz transmission field represented as a
%   homogeneous Rayleigh field plus common-density interface layer potentials.
%
% Main workflow:
%   Reconstruct the source statement, define the Sobolev and modal setting,
%   prove the homogeneous Fourier expansion, derive the direct Green formula,
%   isolate the extra Muller synthesis theorem, and separate field, component,
%   coefficient, and density uniqueness.
%
% Based on:
%   draft/draft.tex, especially Conjecture 1 and the surrounding definitions;
%   Barnett--Greengard (2010), Appendix A; Colton--Kress (1983/2013),
%   Theorem 3.41; Linton (1998); and standard transmission layer-potential
%   theory.
%
% Main changes:
%   The original imprecise uniqueness claim is replaced by an unconditional
%   direct Green representation and a corrected common-density theorem.  The
%   mixed QP/free ansatz leads to a standard full-plane complementary
%   transmission problem, not a quasiperiodic-cylinder spectral condition.
%   The synthesis proof is expanded into an explicit complementary field,
%   four Green extinction formulas, and closed density formulas.
%
% Numerical goal:
%   None. This file contains theory only and no numerical experiment.

\documentclass[11pt]{article}
\usepackage[margin=1in]{geometry}
\usepackage{amsmath,amssymb,amsthm,mathtools,mathrsfs}
\usepackage{booktabs,tabularx}
\usepackage{enumitem}
\usepackage{microtype}
\usepackage{xcolor}
\usepackage{hyperref}
\hypersetup{colorlinks=true,linkcolor=blue!55!black,citecolor=blue!55!black,urlcolor=blue!55!black}
\setlist{itemsep=2pt,topsep=4pt}
\allowdisplaybreaks

\newcommand{\C}{\mathcal C}
\newcommand{\Om}{\Omega}
\newcommand{\Sig}{\Sigma}
\newcommand{\R}{\mathbb R}
\newcommand{\Z}{\mathbb Z}
\newcommand{\ii}{\mathrm i}
\newcommand{\ee}{\mathrm e}
\newcommand{\QP}{\mathrm{QP}}
\newcommand{\Tr}{\operatorname{Tr}}
\newcommand{\Ran}{\operatorname{Ran}}
\newcommand{\Ker}{\operatorname{Ker}}
\newcommand{\diag}{\operatorname{diag}}
\newcommand{\pair}[2]{\left\langle #1,#2\right\rangle}

\theoremstyle{plain}
\newtheorem{theorem}{Theorem}[section]
\newtheorem{proposition}[theorem]{Proposition}
\newtheorem{lemma}[theorem]{Lemma}
\newtheorem{corollary}[theorem]{Corollary}
\theoremstyle{definition}
\newtheorem{definition}[theorem]{Definition}
\theoremstyle{remark}
\newtheorem{remark}[theorem]{Remark}

\title{Conjecture 1: Single-Cell Helmholtz Fields as a Homogeneous Background Plus Interface Potentials}
\author{A theoretical audit for the TEP project}
\date{14 July 2026}

\begin{document}
\maketitle

\begin{abstract}
The TEP draft conjectures that every sufficiently regular, vertically
quasiperiodic Helmholtz transmission field in one cell, with no condition on
the two vertical ports, has a unique representation as a homogeneous
Rayleigh--Bloch field plus quasiperiodic exterior and free-space interior layer
potentials using one common density pair.  The claim is not a valid theorem as
written because its spaces, radiation convention, and meaning of uniqueness are
unspecified.  This report proves an unconditional Green representation with
physical Cauchy data and opposite interior/exterior signs.  It then proves a
corrected version of the draft's common-density representation.  The key
non-elementary step is the standard well-posedness of a swapped-wavenumber
full-plane transmission problem: the complementary exterior field comes from
the nonphysical extension of the free-space interior potential and therefore
satisfies the radial Sommerfeld condition, not a quasiperiodic-cylinder
radiation condition.  Away from Wood thresholds, and after fixing the incident
normalization of the two Rayleigh families, the component fields, Rayleigh
coefficients, and common density pair are unique.  At a Wood threshold the
displayed exponential basis misses the generalized mode \(x\psi_m\).  Physical
guided modes remain intended null fields of the mismatch operator and do not
invalidate the representation.
\end{abstract}

\tableofcontents

\section{Introduction and verdict}

The representation in question is not merely a convenient formula.  It is the
surjectivity and injectivity assertion behind the later block system combining
a Muller interface row with Rayleigh coefficients and lead traces.  If the
representation misses fields, the block system can miss eigenfunctions.  If it
has coordinate nullspaces, a singular discretized matrix need not correspond
to a nonzero physical field.

The conclusion of this investigation is deliberately split into two levels.

\begin{enumerate}
\item An \emph{unconditional direct Green representation} is available under
standard trace assumptions.  Its interface densities are the physical
Dirichlet and Neumann traces, but the interior and exterior interface terms
have opposite signs.
\item The exact \emph{same-density, same-sign Muller ansatz} in the draft is an
indirect representation.  Its completeness is proved through a complementary
(swapped-wavenumber) transmission problem.  Because the ansatz uses a QP kernel
outside and a free-space kernel inside, the complementary exterior field is a
free-space radiating field on \(\R^2\setminus\overline\Om\).  Standard
transmission well-posedness supplies the needed existence and uniqueness; no
additional quasiperiodic-cylinder nonresonance hypothesis is introduced.
\end{enumerate}

Consequently, Conjecture 1 is classified as a \emph{correctable claim}: after
the spaces, non-Wood condition, port normalization, and kernel conventions are
made explicit, the corrected theorem in Section~\ref{sec:corrected} is proved.

\section{The source statement and what it leaves open}
\label{sec:source}

The source is \texttt{draft/draft.tex}, lines 583--622.  With notation retained
from the draft, its exact mathematical content is:

\begin{quote}
Let \(\C=[X_L,X_R]\times[0,d]\) contain an inclusion \(\Om\), and assume
\((k,\beta)\) is not at a Wood anomaly.  Every sufficiently regular field
satisfying
\[
\begin{cases}
\Delta u+k^2u=0&\text{in }\C\setminus\overline\Om,\\
\Delta u+(nk)^2u=0&\text{in }\Om,\\
u(x,d)=\ee^{\ii\beta d}u(x,0),\quad
u_y(x,d)=\ee^{\ii\beta d}u_y(x,0),\\
u,\ u_\nu\text{ continuous on }\partial\Om,
\end{cases}
\]
with no vertical-port boundary condition has a unique representation
\[
u=\begin{cases}
S_{\QP}^{(k)}\sigma+D_{\QP}^{(k)}\tau+u_h[\xi],
&\C\setminus\overline\Om,\\
S^{(nk)}\sigma+D^{(nk)}\tau,&\Om,
\end{cases}
\]
where the Rayleigh series for \(u_h\) converges in an appropriate \(H^1\)
sense.
\end{quote}

The text defines
\begin{equation}
u_h(x,y)=\sum_{m\in\Z}a_m\ee^{\ii\gamma_m(x-X_L)}\psi_m(y)
+\sum_{m\in\Z}b_m\ee^{-\ii\gamma_m(x-X_R)}\psi_m(y),
\label{eq:uh}
\end{equation}
with
\begin{equation}
\beta_m=\beta+\frac{2\pi m}{d},\qquad
\psi_m(y)=d^{-1/2}\ee^{\ii\beta_my},\qquad
\gamma_m^2=k^2-\beta_m^2,\qquad \Im\gamma_m\geq0.
\label{eq:modes}
\end{equation}
For the corrected theorem, the branch convention is strengthened to
\(\Re\gamma_m\geq0\), \(\Im\gamma_m\geq0\), with \(\gamma_m>0\) for a
nonzero propagating order.  The draft's condition \(\Im\gamma_m\geq0\) alone
does not choose between the two real square roots.

Several logically different meanings are hidden in the word ``unique'':
uniqueness of the PDE field, of \(u_h\), of the layer field, of the density
pair, and of the two Rayleigh coefficient sequences.  The PDE field is
obviously not unique because no vertical boundary condition is imposed.  The
only sensible question is uniqueness of coordinates for a \emph{given} field.
The source also leaves the boundary regularity, density spaces, jump signs,
coefficient norm, and exceptional transmission spectrum unspecified.

\section{Geometry, weak solutions, and trace spaces}
\label{sec:spaces}

\subsection{The quasiperiodic cylinder}

It is convenient to identify the horizontal sides of the cell and work on
\[
\mathscr X_\beta=\R\times(\R/d\Z),\qquad
v(x,y+d)=\ee^{\ii\beta d}v(x,y).
\]
Let \(\Sig=\partial\Om\) be of class \(C^2\), let
\(\overline\Om\subset(X_L,X_R)\times(0,d)\), and assume a positive distance
between \(\Sig\) and both vertical ports.  The normal \(\nu\) points from
\(\Om\) into the background.  We take \(k>0\), \(n>0\), and real \(\beta\),
with
\[
k_e=k,\qquad k_i=nk.
\]
The absorbing case \(\Im k>0\) is easier because outgoing uniqueness follows
from an energy identity; the real lossless case contains the spectral issues
of interest.

\subsection{What \texorpdfstring{\(H^1\), \(H^{1/2}\), and \(H^{-1/2}\)}{H1, H1/2, and H-1/2} mean}

The space \(L^2(D)\) consists of square-integrable functions.  A function is in
\(H^1(D)\) if it and its first \emph{weak} derivatives are in \(L^2(D)\).  A
weak derivative is defined by moving a derivative onto a smooth test function
in the integration-by-parts identity; no pointwise differentiability is
required.

On a Lipschitz boundary, the trace theorem gives a bounded map
\[
\gamma:H^1(D)\longrightarrow H^{1/2}(\partial D).
\]
For a Helmholtz solution \(v\in H^1(D)\) with \(\Delta v\in L^2(D)\), Green's
identity defines its normal derivative by
\begin{equation}
\pair{\partial_\nu v}{\gamma\phi}_{\partial D}
=\int_D \nabla v\cdot\nabla\overline\phi-k^2v\overline\phi\,dx.
\label{eq:weaknormal}
\end{equation}
This is a continuous functional on \(H^{1/2}(\partial D)\), hence an element of
its dual \(H^{-1/2}(\partial D)\).  Thus Dirichlet and Neumann traces naturally
live in \(H^{1/2}\) and \(H^{-1/2}\), respectively.  These facts are standard
consequences of the trace and weak Green theories; see McLean
\cite[Chapters 3--4]{McLean2000}.

We use the solution class
\begin{equation}
u_e\in H^1_\beta(\C\setminus\overline\Om),\qquad
u_i\in H^1(\Om),
\label{eq:solutionclass}
\end{equation}
with weak Helmholtz equations and
\begin{equation}
\gamma^+u_e=\gamma^-u_i\in H^{1/2}(\Sig),\qquad
\partial_\nu^+u_e=\partial_\nu^-u_i\in H^{-1/2}(\Sig).
\label{eq:transmission}
\end{equation}
One can first prove all formulas for piecewise \(H^2\) fields and then pass to
\eqref{eq:solutionclass} by continuity.

\subsection{Quasiperiodic Sobolev spaces}

If \(f=\sum_m f_m\psi_m\), define
\begin{equation}
\|f\|_{H^s_\beta(0,d)}^2
=\sum_{m\in\Z}(1+|\beta_m|^2)^s|f_m|^2.
\label{eq:Hs}
\end{equation}
This is the periodic Fourier definition after factoring out
\(\ee^{\ii\beta y}\).  It also makes the duality
\((H^{1/2}_\beta)'=H^{-1/2}_\beta\) transparent: Cauchy--Schwarz pairs the
weights \((1+|\beta_m|^2)^{1/2}\) and its reciprocal.

\section{Homogeneous Rayleigh representation}
\label{sec:rayleigh}

\subsection{Modewise derivation}

Let \(h\in H^1_\beta((A,B)\times(0,d))\) solve
\((\Delta+k^2)h=0\) weakly.  Expanding in the orthonormal basis \(\psi_m\),
\[
h(x,y)=\sum_m h_m(x)\psi_m(y),
\]
and testing the weak equation with \(\varphi(x)\overline{\psi_m(y)}\) gives
\begin{equation}
h_m''+(k^2-\beta_m^2)h_m=0
\quad\text{in }\mathcal D'(A,B).
\label{eq:ode}
\end{equation}
Distributional solutions of this constant-coefficient ODE are smooth.  If
\(\gamma_m\neq0\),
\begin{equation}
h_m(x)=a_m\ee^{\ii\gamma_m(x-X_L)}
+b_m\ee^{-\ii\gamma_m(x-X_R)}.
\label{eq:modepair}
\end{equation}
The two functions have Wronskian \(-2\ii\gamma_m\) times a nonzero phase and
are therefore linearly independent.  This proves modewise uniqueness away
from thresholds.

At \(\gamma_m=0\), equation \eqref{eq:ode} becomes \(h_m''=0\), whose solution
space is
\begin{equation}
h_m(x)=A_m+B_mx.
\label{eq:threshold}
\end{equation}
The two exponentials in \eqref{eq:modepair} both become 1 and cannot represent
the second solution.  This is an explicit, elementary failure of the displayed
draft coordinates at a Wood anomaly.

\subsection{Coefficient norm and convergence}

For large \(|m|\), write \(\gamma_m=\ii\alpha_m\), \(\alpha_m>0\).  The two
anchored factors are
\[
\ee^{-\alpha_m(x-X_L)},\qquad
\ee^{-\alpha_m(X_R-x)}.
\]
They decay into the cell from opposite ports.  Direct integration of the
function, \(x\)-derivative, and \(y\)-derivative terms gives, up to constants
depending only on the cell width,
\begin{equation}
\|h\|_{H^1(\C)}^2\asymp
\sum_m(1+|\beta_m|^2)^{1/2}(|a_m|^2+|b_m|^2),
\label{eq:coefnorm}
\end{equation}
apart from finitely many propagating modes.  The cross terms contain
\(\ee^{-\alpha_m(X_R-X_L)}\) and are harmless.  Thus the appropriate space is
\[
\ell^2_{1/2,\beta}\times\ell^2_{1/2,\beta}.
\]
Equation \eqref{eq:coefnorm} proves convergence in \(H^1\), convergence of
Dirichlet traces in \(H^{1/2}\), and convergence of normal traces in
\(H^{-1/2}\).

\begin{proposition}[Empty-cell check]
\label{prop:empty}
If \(\Om=\varnothing\) and \(\gamma_m\neq0\) for all \(m\), every weak
\(H^1_\beta\) homogeneous cell solution has a unique representation
\eqref{eq:uh} with coefficients in
\(\ell^2_{1/2,\beta}\times\ell^2_{1/2,\beta}\).  At a threshold, the formula
must be enlarged by \eqref{eq:threshold}.
\end{proposition}

\begin{proof}
Equations \eqref{eq:ode}--\eqref{eq:coefnorm} give existence, convergence, and
modewise uniqueness.  Orthogonality of \(\psi_m\) prevents cancellation
between different orders.
\end{proof}

\subsection{Incoming normalization}

Near the left port, write the local field as
\[
u_e=\sum_m(A_m^L\ee^{\ii\gamma_m(x-X_L)}
+B_m^L\ee^{-\ii\gamma_m(x-X_L)})\psi_m.
\]
Near the right port use coefficients \(A_m^R,B_m^R\) anchored at \(X_R\).
For a non-Wood order, Dirichlet and \(x\)-derivative Fourier coefficients
\((f_m,g_m)\) determine
\begin{equation}
A_m=\tfrac12\left(f_m+\frac{g_m}{\ii\gamma_m}\right),\qquad
B_m=\tfrac12\left(f_m-\frac{g_m}{\ii\gamma_m}\right).
\label{eq:portcoeff}
\end{equation}
The trace weights show that \(f\in H^{1/2}\) and \(g\in H^{-1/2}\) imply the
half-order coefficient bound.

We define the background field by the \emph{incoming} coefficients
\begin{equation}
u_h=\sum_m A_m^L\ee^{\ii\gamma_m(x-X_L)}\psi_m
+\sum_m B_m^R\ee^{-\ii\gamma_m(x-X_R)}\psi_m.
\label{eq:incomingh}
\end{equation}
Then \(w_e=u_e-u_h\) is outgoing at both ends: it contains only the left-going
factor on the left and the right-going factor on the right.  For evanescent
orders ``outgoing'' means exponential decay.  This normalization is essential;
without saying how \(u_h\) is tied to the two ports, the original uniqueness
claim has no precise content.

\section{Quasiperiodic Green function and layer potentials}
\label{sec:layers}

Let
\[
\Phi_\kappa(x,z)=\frac{\ii}{4}H_0^{(1)}(\kappa|x-z|),
\]
and define
\begin{equation}
S_\kappa\sigma(x)=\int_\Sig\Phi_\kappa(x,z)\sigma(z)\,ds_z,
\quad
D_\kappa\tau(x)=\int_\Sig\partial_{\nu_z}\Phi_\kappa(x,z)\tau(z)\,ds_z.
\label{eq:potentials}
\end{equation}
The exterior operators replace \(\Phi_k\) by the outgoing
\(y\)-quasiperiodic Green function.  Its spectral form, in the draft's sign
convention, is
\begin{equation}
G_{\QP}^{(k)}(x,y;x',y')
=\frac{\ii}{2d}\sum_{m\in\Z}\frac{1}{\gamma_m}
\ee^{\ii\beta_m(y-y')}\ee^{\ii\gamma_m|x-x'|}.
\label{eq:qpgreen}
\end{equation}
Linton derives the equivalent eigenfunction expansion and its outgoing
asymptotics in equations (2.13)--(2.17) of \cite{Linton1998}.  Formula
\eqref{eq:qpgreen} is undefined when a \(\gamma_m\) vanishes.

For the normal pointing out of \(\Om\), write \(+\) for the trace from the
exterior and \(-\) for the trace from the interior, and choose the jump
convention
\begin{align}
\gamma^\pm D\tau&=K\tau\pm\tfrac12\tau,
&\gamma^\pm S\sigma&=S\sigma,\label{eq:jumps1}\\
\partial_\nu^\pm S\sigma&=K'\sigma\mp\tfrac12\sigma,
&\partial_\nu^\pm D\tau&=T\tau.\label{eq:jumps2}
\end{align}
Changing the sign of the fundamental solution or normal changes all displayed
half signs together; it does not change the conclusions.  The mapping spaces
are
\[
\tau\in H^{1/2}(\Sig),\qquad \sigma\in H^{-1/2}(\Sig).
\]
The trace and mapping statements in these spaces follow from the Calderon
theory summarized in \cite[Sections 3--4]{DominguezLyonTurc2016}.

\section{What direct Green's identity actually proves}
\label{sec:green}

Let
\[
p=\gamma u\in H^{1/2}(\Sig),\qquad
q=\partial_\nu u\in H^{-1/2}(\Sig).
\]
Apply Green's second identity first in \(\Om\) and then in
\(\C\setminus\overline\Om\).  In the latter domain, the outward normal on the
inner boundary is \(-\nu\).  The top and bottom integrals cancel: the field has
phase \(\ee^{\ii\beta d}\), while the Green kernel as a function of the source
variable has the inverse phase; the normal directions are opposite.  This is
the one-period analogue of Lemma 10 in \cite[Appendix A]{BarnettGreengard2010}.
The two vertical-boundary integrals remain.

Denote their sum by \(h_V\).  For targets in the open cell, its sources lie on
the vertical boundary, so
\((\Delta+k^2)h_V=0\) and it is \(y\)-quasiperiodic.  Section
\ref{sec:rayleigh} therefore gives it a two-sided Rayleigh expansion in the
cell.  Careful use of the inner-boundary normal gives
\begin{theorem}[Direct Green representation]
\label{thm:directgreen}
Under the geometry and weak-solution assumptions of Section~\ref{sec:spaces},
and away from Wood thresholds,
\begin{align}
u_e&=h_V+D_{\QP}^{(k_e)}p-S_{\QP}^{(k_e)}q
&&\text{in }\C\setminus\overline\Om,\label{eq:directext}\\
u_i&=S^{(k_i)}q-D^{(k_i)}p
&&\text{in }\Om.\label{eq:directint}
\end{align}
Here \(h_V\) is the retained vertical-boundary Cauchy integral and has an
\(H^1\)-convergent Rayleigh expansion.
\end{theorem}

\begin{proof}
For piecewise \(H^2\), excise a small circle around the target, apply Green's
second identity, and let its radius tend to zero.  The small-circle term gives
the field value.  The horizontal terms cancel as described above.  On
\(\Sig\), the background-domain normal is \(-\nu\), which changes both the
normal derivative of the field and the source-normal derivative of the Green
kernel and yields \(+Dp-Sq\).  In the inclusion the ordinary interior formula
yields \(+Sq-Dp\).  For piecewise \(H^1\), interpret normal derivatives by
\eqref{eq:weaknormal}; boundedness of trace and potential maps extends the
identity by density.
\end{proof}

The signs in \eqref{eq:directext}--\eqref{eq:directint} are opposite.  Hence
direct Green's identity does \emph{not} prove the draft's formula with the same
\((\sigma,\tau)\) and the same signs on both sides.  That formula is a Muller
ansatz and needs an additional completeness theorem.

\section{The common-density Muller synthesis}
\label{sec:muller}

Define the synthesis map
\begin{equation}
\mathcal J(\tau,\sigma)=
\left(
D_{\QP}^{(k_e)}\tau+S_{\QP}^{(k_e)}\sigma,
D^{(k_i)}\tau+S^{(k_i)}\sigma
\right).
\label{eq:J}
\end{equation}
Its exterior member is outgoing because of \eqref{eq:qpgreen}.  Using
\eqref{eq:jumps1}--\eqref{eq:jumps2} and the draft vector
\(\eta=(\tau,-\sigma)^T\), its Cauchy mismatch is
\begin{equation}
A_{\QP}\eta=
\begin{bmatrix}
I+K_e-K_i&S_i-S_e\\
T_e-T_i&I+K_i'-K_e'
\end{bmatrix}\eta.
\label{eq:A}
\end{equation}
This recovers the operator printed after Conjecture 1.  The leading
singularities cancel in each wavenumber difference.  On a \(C^2\) curve,
\(A_{\QP}\) is identity plus compact on
\begin{equation}
\mathcal X=H^{1/2}(\Sig)\times H^{-1/2}(\Sig).
\label{eq:X}
\end{equation}
Barnett and Greengard make the same observation after their equation (20)
\cite[Section 3]{BarnettGreengard2010}.  Compactness of corresponding
wavenumber differences in Sobolev trace spaces is recorded in
\cite[Section 4.1]{DominguezLyonTurc2016}.

\subsection{The complementary problem}

The crucial object is not an interior Dirichlet or Neumann eigenproblem.  It is
the \emph{swapped} transmission problem obtained by extending the two
potentials into their nonphysical sides:
\begin{equation}
\begin{cases}
(\Delta+k_e^2)v_i=0&\text{in }\Om,\\
(\Delta+k_i^2)v_e=0&\text{in }\R^2\setminus\overline\Om,\\
v_e-v_i=F,\quad
\partial_\nu v_e-\partial_\nu v_i=G&\text{on }\Sig,\\
\displaystyle
\partial_r v_e-\ii k_i v_e=o(r^{-1/2})
&\text{as }r=|x|\to\infty,
\end{cases}
\label{eq:complement}
\end{equation}
The wavenumbers have exchanged physical locations.  The geometry is determined
by the kernels being continued, not by the physical exterior geometry.  The
nonphysical continuation of the QP exterior potential gives \(v_i\) inside
\(\Om\), while the nonphysical continuation of the free-space interior
potential gives \(v_e\) on the full plane outside \(\Om\).  Thus \(v_e\)
satisfies the radial Sommerfeld condition.  This is exactly the
``complementary field'' device in Appendix A of
\cite{BarnettGreengard2010}.

\begin{remark}[Relation to Hiptmair--Moiola--Spence (2022)]
Definition 1.4 of \cite{HiptmairMoiolaSpence2022} distinguishes the physical
solution operator \(S_{io}\) from the swapped-parameter unphysical operator
\(S_{oi}\), and Theorem 1.10 expresses standard transmission-BIE inverses in
terms of both.  This supports the structural distinction used here.  Their
paper concerns full-plane direct BIEs; the present common-density calculation
uses the mixed QP/free indirect ansatz of \cite[Appendix A]{BarnettGreengard2010}.
\end{remark}

\begin{definition}[Well-posed complementary problem]
\label{def:complement-solvability}
The complementary problem \eqref{eq:complement} is called well posed if, for
every
\[
(F,G)\in H^{1/2}(\Sig)\times H^{-1/2}(\Sig),
\]
there exists exactly one pair
\[
v_i\in H^1(\Om),\qquad
v_e\in H^1_{\rm loc}(\R^2\setminus\overline\Om)
\]
satisfying \eqref{eq:complement} and the Sommerfeld radiation condition.
\end{definition}

For a bounded \(C^2\) inclusion and positive real \(k_e,k_i\), this is the
standard free-space Helmholtz transmission problem.  Colton and Kress state
its classical existence and uniqueness as ``The transmission problem has a
unique solution'' in Theorem 3.41 of \cite[Theorem 3.41,
p.~101]{ColtonKress1983}; their preceding Theorem 3.40 proves uniqueness.
Barnett and Greengard invoke exactly those two results in their Appendix A
\cite[pp.~22--24]{BarnettGreengard2010}.  Colton and Kress formulate the
argument for classical traces on a \(C^2\) surface in \(\R^3\).  The same
layer-potential/Fredholm proof, with the two-dimensional outgoing fundamental
solution and the Sobolev mapping properties in Section~\ref{sec:layers}, gives
Definition~\ref{def:complement-solvability}.  Thus this is standard
full-plane transmission well-posedness, not an additional periodic-waveguide
spectral assumption.

The complementary field can itself be written as layer potentials.  For any
\((\tau,\sigma)\in\mathcal X\), define its physical fields
\begin{equation}
P_e=D_{\QP}^{(k_e)}\tau+S_{\QP}^{(k_e)}\sigma
\quad\hbox{in }\C\setminus\overline\Om,
\qquad
P_i=D^{(k_i)}\tau+S^{(k_i)}\sigma
\quad\hbox{in }\Om,
\label{eq:physical-potentials}
\end{equation}
and continue the same potentials to their nonphysical sides with one sign
change:
\begin{equation}
v^{\tau,\sigma}(x)=
\begin{cases}
v_i(x):=D_{\QP}^{(k_e)}\tau(x)+S_{\QP}^{(k_e)}\sigma(x),
&x\in\Om,\\
v_e(x):=-D^{(k_i)}\tau(x)-S^{(k_i)}\sigma(x),
&x\in\R^2\setminus\overline\Om.
\end{cases}
\label{eq:complementary-layer}
\end{equation}
This is the present sign convention's version of equation (A.7) of
\cite{BarnettGreengard2010}.  The QP potential has merely been evaluated on
the other side of \(\Sig\), while the free-space potential in the second line
is radiating because it is a finite-interface outgoing layer potential.

\begin{lemma}[Explicit common-density synthesis]
\label{lem:synthesis}
Assume non-Wood parameters and the geometry and positive-material hypotheses
above.  Let
\[
w_e\in H^1_{\rm loc}(\C\setminus\overline\Om),
\qquad w_i\in H^1(\Om),
\]
where \(w_e\) is \(y\)-quasiperiodic, solves
\((\Delta+k_e^2)w_e=0\), and is outgoing at both vertical ends, while
\((\Delta+k_i^2)w_i=0\) in \(\Om\).  No transmission condition between
\(w_e\) and \(w_i\) is assumed.  Then there is exactly one
\((\tau,\sigma)\in\mathcal X\) such that
\begin{equation}
w_e=D_{\QP}^{(k_e)}\tau+S_{\QP}^{(k_e)}\sigma,
\qquad
w_i=D^{(k_i)}\tau+S^{(k_i)}\sigma
\label{eq:synthesis-target}
\end{equation}
in their respective domains.
\end{lemma}

\begin{proof}
We prove uniqueness and existence separately.

\emph{1. Uniqueness of the densities.}
Suppose first that both physical fields in
\eqref{eq:physical-potentials} vanish.  Define \((v_i,v_e)\) by
\eqref{eq:complementary-layer}.  Since \(P_e\) is zero on its exterior side,
the jumps \eqref{eq:jumps1}--\eqref{eq:jumps2} give
\begin{equation}
\gamma v_i=-\tau,\qquad \partial_\nu v_i=\sigma.
\label{eq:kernel-traces-i}
\end{equation}
Since \(P_i\) is zero on its interior side and \(v_e=-P_i\) on the exterior
side, the same jumps give
\begin{equation}
\gamma v_e=-\tau,\qquad \partial_\nu v_e=\sigma.
\label{eq:kernel-traces-e}
\end{equation}
Thus \((v_i,v_e)\) satisfies \eqref{eq:complement} with \(F=G=0\).  The
uniqueness part of Definition~\ref{def:complement-solvability} yields
\(v_i=v_e=0\).  Equations \eqref{eq:kernel-traces-i}--
\eqref{eq:kernel-traces-e} then give \(\tau=\sigma=0\).  Hence
\begin{equation}
\Ker\mathcal J=\{0\}.
\label{eq:Jinjective}
\end{equation}

\emph{2. Construction of the densities.}
For the prescribed fields in the lemma, introduce the four Cauchy traces
\begin{equation}
p_e=\gamma^+w_e,\quad q_e=\partial_\nu^+w_e,
\qquad
p_i=\gamma^-w_i,\quad q_i=\partial_\nu^-w_i.
\label{eq:w-traces}
\end{equation}
Use Definition~\ref{def:complement-solvability} to solve
\eqref{eq:complement} with
\begin{equation}
F=-(p_e+p_i),\qquad G=-(q_e+q_i).
\label{eq:complement-data-from-w}
\end{equation}
Write
\[
p_{v,i}=\gamma v_i,\quad q_{v,i}=\partial_\nu v_i,
\qquad
p_{v,e}=\gamma v_e,\quad q_{v,e}=\partial_\nu v_e.
\]
The two jump equations in \eqref{eq:complement} are precisely
\begin{align}
p_{v,e}-p_{v,i}&=-(p_e+p_i),\label{eq:vjump-p}\\
q_{v,e}-q_{v,i}&=-(q_e+q_i).\label{eq:vjump-q}
\end{align}
Therefore the following two expressions on each line agree, and we define
\begin{align}
\tau&:=p_e-p_{v,i}=-p_{v,e}-p_i,\label{eq:density-tau}\\
\sigma&:=q_{v,i}-q_e=q_{v,e}+q_i.\label{eq:density-sigma}
\end{align}

It remains to verify that these densities produce the prescribed fields.  We
first prove the two quasiperiodic Green formulas separately, in the style of
Lemmas 9--10 of \cite[Appendix A]{BarnettGreengard2010}.

\begin{lemma}[Interior quasiperiodic Green formula]
\label{lem:qp-green-interior}
Let \(v_i\) satisfy \((\Delta+k_e^2)v_i=0\) in \(\Om\), and let
\(p_{v,i}=\gamma v_i\), \(q_{v,i}=\partial_\nu v_i\).  Then
\begin{equation}
D_{\QP}^{(k_e)}p_{v,i}-S_{\QP}^{(k_e)}q_{v,i}
=\begin{cases}
-v_i,&\text{in }\Om,\\
0,&\text{in }\C\setminus\overline\Om.
\end{cases}
\label{eq:green-vi}
\end{equation}
\end{lemma}

\begin{proof}
Put \(e_y=(0,1)\), \(\alpha=\ee^{\ii\beta d}\), and
\(\Om_j=\Om+jd e_y\), \(\Sig_j=\partial\Om_j\).  The one-period
quasiperiodic Green function is the appropriately converged lattice sum
\[
G_{\QP}^{(k_e)}(x,z)
=\sum_{j\in\Z}\alpha^j\Phi_{k_e}(x,z+jd e_y).
\]
For clarity, first truncate this sum to \(|j|\leq N\).  On \(\Sig_j\), copy
the Cauchy data by
\[
p_{v,i}^{(j)}(z+jd e_y)=p_{v,i}(z),\qquad
q_{v,i}^{(j)}(z+jd e_y)=q_{v,i}(z).
\]
After the change of variable \(z_j=z+jd e_y\), the contribution of the
\(j\)-th image is
\[
\begin{aligned}
I_j(x)
&=\int_{\Sig_j}
 \left[\partial_{\nu_{z_j}}\Phi_{k_e}(x,z_j)
             p_{v,i}^{(j)}(z_j)
       -\Phi_{k_e}(x,z_j)q_{v,i}^{(j)}(z_j)\right]ds_{z_j}\\
&=\begin{cases}
-v_i(x-jd e_y),&x\in\Om_j,\\
0,&x\notin\overline\Om_j.
\end{cases}
\end{aligned}
\]
This is exactly the ordinary interior Green formula together with its
extinction case.  Since \(\overline\Om\subset\R\times(0,d)\), a target in the
open reference cell belongs to at most one translated inclusion.  Hence
\[
\sum_{|j|\leq N}\alpha^j I_j(x)
=\begin{cases}-v_i(x),&x\in\Om,\\0,&x\in\C\setminus\overline\Om,
\end{cases}
\]
for every sufficiently large \(N\).  Passing to the non-Wood quasiperiodic
Green-function limit proves \eqref{eq:green-vi}.  This is the one-period
version of Lemma 9 of \cite[Appendix A]{BarnettGreengard2010}.
\end{proof}

\begin{lemma}[Outgoing exterior quasiperiodic Green formula]
\label{lem:qp-green-exterior}
Let \(w_e\) satisfy \((\Delta+k_e^2)w_e=0\) in
\(\C\setminus\overline\Om\), be \(y\)-quasiperiodic with phase
\(\alpha=\ee^{\ii\beta d}\), and be outgoing at both vertical ends.  With
\(p_e=\gamma^+w_e\) and \(q_e=\partial_\nu^+w_e\),
\begin{equation}
D_{\QP}^{(k_e)}p_e-S_{\QP}^{(k_e)}q_e
=\begin{cases}
0,&\text{in }\Om,\\
w_e,&\text{in }\C\setminus\overline\Om.
\end{cases}
\label{eq:green-we}
\end{equation}
\end{lemma}

\begin{proof}
Fix a target \(x=(x_1,x_2)\) and write the source point as \(z=(s,t)\).
Choose \(L<X_L\) and \(R>X_R\), with \(L<x_1<R\), and set
\[
\mathcal D_{L,R}=((L,R)\times(0,d))\setminus\overline\Om.
\]
If \(x\in\mathcal D_{L,R}\), remove a disk \(B_\varepsilon(x)\); if
\(x\in\Om\), no disk is removed.  Apply Green's second identity to
\(w_e(z)\) and \(G_{\QP}^{(k_e)}(x,z)\).  The boundary integrand is
\[
\mathcal B(z)
=w_e(z)\partial_{n_z}G_{\QP}^{(k_e)}(x,z)
 -G_{\QP}^{(k_e)}(x,z)\partial_{n_z}w_e(z).
\]

We first dispose of the horizontal sides.  As a function of its source,
the Green function has the inverse quasiperiodicity
\[
\begin{aligned}
w_e(s,d)&=\alpha w_e(s,0),
&\partial_t w_e(s,d)&=\alpha\partial_t w_e(s,0),\\
G_{\QP}(x;s,d)&=\alpha^{-1}G_{\QP}(x;s,0),
&\partial_tG_{\QP}(x;s,d)&=\alpha^{-1}\partial_tG_{\QP}(x;s,0).
\end{aligned}
\]
Because the outward normals are \(+e_y\) at \(t=d\) and \(-e_y\) at
\(t=0\), the two integrals are
\[
\begin{aligned}
I_T
&=\int_L^R\left[
 \alpha w_e(s,0)\,\alpha^{-1}\partial_tG_{\QP}(x;s,0)
 -\alpha^{-1}G_{\QP}(x;s,0)\,\alpha\partial_tw_e(s,0)
 \right]ds,\\
I_B
&=-\int_L^R\left[
 w_e(s,0)\partial_tG_{\QP}(x;s,0)
 -G_{\QP}(x;s,0)\partial_tw_e(s,0)
 \right]ds.
\end{aligned}
\]
Thus \(I_T+I_B=0\), exactly as in Barnett's Lemma 10.

It remains to examine the two vertical sides.  Define the anti-quasiperiodic
dual modes
\[
\widetilde\psi_m(t)=d^{-1/2}\ee^{-\ii\beta_m t},\qquad
\int_0^d\psi_m(t)\widetilde\psi_n(t)\,dt=\delta_{mn}.
\]
The outgoing hypothesis gives, on the right and left homogeneous ports,
\[
\begin{aligned}
w_e(s,t)&=\sum_m A_m^R
 \ee^{\ii\gamma_m(s-X_R)}\psi_m(t),&&s\geq X_R,\\
w_e(s,t)&=\sum_m B_m^L
 \ee^{-\ii\gamma_m(s-X_L)}\psi_m(t),&&s\leq X_L.
\end{aligned}
\]
The spectral Green formula \eqref{eq:qpgreen}, now viewed as a function of
the source \(z=(s,t)\), gives
\[
\begin{aligned}
G_{\QP}(x;s,t)
 &=\sum_m c_m^R(x)\ee^{\ii\gamma_m(s-X_R)}
   \widetilde\psi_m(t),&&s>x_1,\\
G_{\QP}(x;s,t)
 &=\sum_m c_m^L(x)\ee^{-\ii\gamma_m(s-X_L)}
   \widetilde\psi_m(t),&&s<x_1,
\end{aligned}
\]
where, explicitly,
\[
c_m^R(x)=\frac{\ii}{2\sqrt d\,\gamma_m}
 \ee^{\ii\beta_mx_2}\ee^{\ii\gamma_m(X_R-x_1)},\qquad
c_m^L(x)=\frac{\ii}{2\sqrt d\,\gamma_m}
 \ee^{\ii\beta_mx_2}\ee^{\ii\gamma_m(x_1-X_L)}.
\]
On the right side \(s=R\), the outward derivative is \(\partial_{n_z}=\partial_s\).
Using the two right-going expansions and the dual-mode identity,
\[
\begin{aligned}
I_R
&=\int_0^d\left(w_e\partial_sG_{\QP}
                 -G_{\QP}\partial_sw_e\right)(x;R,t)\,dt\\
&=\sum_{m,n}A_m^R c_n^R(x)
 \ee^{\ii(\gamma_m+\gamma_n)(R-X_R)}
 \ii(\gamma_n-\gamma_m)
 \int_0^d\psi_m(t)\widetilde\psi_n(t)\,dt\\
&=\sum_m A_m^R c_m^R(x)
 \ee^{2\ii\gamma_m(R-X_R)}\ii(\gamma_m-\gamma_m)=0.
\end{aligned}
\]
On the left side \(s=L\), the outward derivative is
\(\partial_{n_z}=-\partial_s\).  Both left-going factors then have the same
outward logarithmic derivative \(+\ii\gamma_m\), so the identical calculation
gives
\[
\begin{aligned}
I_L
&=\int_0^d\left(w_e\partial_{n_z}G_{\QP}
                 -G_{\QP}\partial_{n_z}w_e\right)(x;L,t)\,dt\\
&=\sum_{m,n}B_m^L c_n^L(x)
 \ee^{-\ii(\gamma_m+\gamma_n)(L-X_L)}
 \ii(\gamma_n-\gamma_m)\delta_{mn}=0.
\end{aligned}
\]
These equalities hold for propagating orders as well as evanescent orders:
decay is not used for a real \(\gamma_m\); the Wronskian vanishes because
\(w_e\) and the Green function contain the same outgoing exponential.  An
incoming exponential would have the opposite logarithmic derivative and
would leave a nonzero term proportional to \(2\ii\gamma_m\).

Finally, on \(\Sig\) the outward normal of \(\mathcal D_{L,R}\) is
\(-\nu\), hence its boundary contribution is
\[
-\int_\Sig\left[p_e(z)\partial_{\nu_z}G_{\QP}(x,z)
                 -G_{\QP}(x,z)q_e(z)\right]ds_z
=-\left(D_{\QP}^{(k_e)}p_e-S_{\QP}^{(k_e)}q_e\right)(x).
\]
If \(x\in\C\setminus\overline\Om\), the small-circle integral tends to
\(+w_e(x)\), since \((\Delta+k_e^2)G_{\QP}(x,\cdot)=-\delta_x\).  All four
outer-side integrals have vanished, so Green's identity becomes
\[
0=-\left(D_{\QP}^{(k_e)}p_e-S_{\QP}^{(k_e)}q_e\right)(x)+w_e(x).
\]
If \(x\in\Om\), there is no small circle and the same calculation gives the
potential equal to zero.  This proves \eqref{eq:green-we}.  For weak fields,
apply the calculation first to finite Rayleigh sums and smooth data, then pass
to the trace-space limit.
\end{proof}

At wavenumber \(k_i\), the two free-space formulas are
\begin{align}
D^{(k_i)}p_{v,e}-S^{(k_i)}q_{v,e}
&=\begin{cases}0,&\text{in }\Om,\\v_e,&\text{in }\R^2\setminus\overline\Om,
\end{cases}\label{eq:green-ve}\\
D^{(k_i)}p_i-S^{(k_i)}q_i
&=\begin{cases}-w_i,&\text{in }\Om,\\0,&\text{in }\R^2\setminus\overline\Om.
\end{cases}\label{eq:green-wi}
\end{align}
Equations \eqref{eq:green-ve}--\eqref{eq:green-wi} are the ordinary exterior
and interior Green formulas.

Now use the first expressions in \eqref{eq:density-tau}--
\eqref{eq:density-sigma}.  In the physical exterior,
\begin{equation}
\begin{aligned}
D_{\QP}^{(k_e)}\tau+S_{\QP}^{(k_e)}\sigma
&=\bigl(D_{\QP}^{(k_e)}p_e-S_{\QP}^{(k_e)}q_e\bigr)
 -\bigl(D_{\QP}^{(k_e)}p_{v,i}-S_{\QP}^{(k_e)}q_{v,i}\bigr)\\
&=w_e-0=w_e.
\end{aligned}
\label{eq:verify-we}
\end{equation}
Use instead the second expressions in
\eqref{eq:density-tau}--\eqref{eq:density-sigma}.  In \(\Om\),
\begin{equation}
\begin{aligned}
D^{(k_i)}\tau+S^{(k_i)}\sigma
&=-\bigl(D^{(k_i)}p_{v,e}-S^{(k_i)}q_{v,e}\bigr)
 -\bigl(D^{(k_i)}p_i-S^{(k_i)}q_i\bigr)\\
&=-0-(-w_i)=w_i.
\end{aligned}
\label{eq:verify-wi}
\end{equation}
This proves existence.  If two density pairs produced the same
\((w_e,w_i)\), their difference would lie in \(\Ker\mathcal J\), which is
zero by \eqref{eq:Jinjective}.  This proves uniqueness.  The calculation is
the full analogue, with the present signs and arbitrary interface mismatch,
of equations (A.7)--(A.14) of \cite{BarnettGreengard2010}.
\end{proof}

\begin{remark}
If both physical layer representations used QP Green functions, their
nonphysical continuations would instead produce a periodic complementary
problem, where spurious guided modes could require a separate spectral
hypothesis.  That is not the ansatz in Conjecture 1: its QP exterior/free-space
interior choice is precisely what makes \eqref{eq:complement} a standard
nonperiodic problem.  Barnett and Greengard emphasize this distinction at the
end of their Appendix A.
\end{remark}

\section{Existence and uniqueness of the corrected representation}
\label{sec:corrected}

\begin{theorem}[Corrected Conjecture 1]
\label{thm:corrected}
Assume:
\begin{enumerate}[label=(\roman*)]
\item \(\Sig\) is \(C^2\), \(\overline\Om\) is strictly inside the cell, and
there are homogeneous background collars at both vertical ports;
\item \(k>0\), \(n>0\), and \(\beta\in\R\), with ordinary TM continuity
\eqref{eq:transmission};
\item \(u\) is a piecewise \(H^1\) weak solution in the sense of
\eqref{eq:solutionclass};
\item \(\gamma_m(k,\beta)\neq0\) for every \(m\).
\end{enumerate}
Define \(u_h\) by the incoming port coefficients in
\eqref{eq:incomingh}.  Then there exists a unique density pair
\((\tau,\sigma)\in H^{1/2}(\Sig)\times H^{-1/2}(\Sig)\) such that
\begin{equation}
u=\begin{cases}
u_h+D_{\QP}^{(k_e)}\tau+S_{\QP}^{(k_e)}\sigma,
&\C\setminus\overline\Om,\\
D^{(k_i)}\tau+S^{(k_i)}\sigma,&\Om.
\end{cases}
\label{eq:correctedrep}
\end{equation}
The Rayleigh coefficients of \(u_h\) belong to
\(\ell^2_{1/2,\beta}\times\ell^2_{1/2,\beta}\) and are unique.  Hence the
background field, layer field, and densities are all unique for the given
\(u\).
\end{theorem}

\begin{proof}
We separate the port normalization, the layer construction, and uniqueness.

\emph{1. Extract the incoming background.}
In the left homogeneous collar, the exterior field has the unique expansion
\begin{equation}
u_e=\sum_m\left(
A_m^L\ee^{\ii\gamma_m(x-X_L)}
+B_m^L\ee^{-\ii\gamma_m(x-X_L)}
\right)\psi_m,
\label{eq:theorem-left-expansion}
\end{equation}
and in the right collar it has the analogous coefficients
\(A_m^R,B_m^R\), anchored at \(X_R\).  Since \(\gamma_m\ne0\), formula
\eqref{eq:portcoeff} recovers all four coefficient sequences from the Cauchy
data at the two ports.  Define the incident background by retaining exactly
the right-going coefficients at the left port and the left-going coefficients
at the right port:
\begin{equation}
u_h=\sum_m A_m^L\ee^{\ii\gamma_m(x-X_L)}\psi_m
+\sum_m B_m^R\ee^{-\ii\gamma_m(x-X_R)}\psi_m.
\label{eq:theorem-uh}
\end{equation}
The estimate \eqref{eq:coefnorm} gives convergence in \(H^1(\C)\).

\emph{2. The remaining exterior field is outgoing.}
Put
\begin{equation}
w_e:=u_e-u_h\quad\text{in }\C\setminus\overline\Om,
\qquad
w_i:=u_i\quad\text{in }\Om.
\label{eq:theorem-w}
\end{equation}
Near the two ports, cancellation of the incoming terms in
\eqref{eq:theorem-left-expansion} and \eqref{eq:theorem-uh} gives
\begin{align}
w_e&=\sum_m B_m^L\ee^{-\ii\gamma_m(x-X_L)}\psi_m
&&\text{near }X_L,\label{eq:w-left-out}\\
w_e&=\sum_m A_m^R\ee^{\ii\gamma_m(x-X_R)}\psi_m
&&\text{near }X_R.\label{eq:w-right-out}
\end{align}
Thus \(w_e\) is outgoing at both ends, including exponential decay for every
evanescent order.  The pair \((w_e,w_i)\) now satisfies every hypothesis of
Lemma~\ref{lem:synthesis}; no interface continuity is needed for that lemma.

\begin{remark}[Why the background is subtracted only from the exterior field]
The background field satisfies
\((\Delta+k_e^2)u_h=0\), whereas the physical interior field satisfies
\((\Delta+k_i^2)u_i=0\).  If one instead defined
\(\widetilde w_i:=u_i-u_h\) in \(\Om\), then
\begin{equation}
  (\Delta+k_i^2)\widetilde w_i
  =(k_e^2-k_i^2)u_h.
  \label{eq:interior-background-subtraction}
\end{equation}
Except when $k_e=k_i$ or $u_h|_\Om=0$, this is a nonzero volume source.
Thus \(\widetilde w_i\) is not a homogeneous interior Helmholtz field and
cannot be represented by the present interface layer potentials alone; a
volume potential would have to be added.  This is why
\eqref{eq:theorem-w} leaves $w_i=u_i$.
\end{remark}

\emph{3. Construct the densities.}
For complete explicitness, let \((p_e,q_e,p_i,q_i)\) be the traces of
\((w_e,w_i)\) in \eqref{eq:w-traces}.  Solve the standard full-plane
complementary problem with
\begin{equation}
v_e-v_i=-(p_e+p_i),\qquad
\partial_\nu v_e-\partial_\nu v_i=-(q_e+q_i).
\label{eq:theorem-complement-data}
\end{equation}
Definition~\ref{def:complement-solvability}, justified by Colton--Kress
Theorem 3.41, gives one and only one \((v_i,v_e)\).  Define
\begin{equation}
\tau=p_e-p_{v,i}=-p_{v,e}-p_i,\qquad
\sigma=q_{v,i}-q_e=q_{v,e}+q_i.
\label{eq:theorem-densities}
\end{equation}
The equalities on each side follow directly from
\eqref{eq:theorem-complement-data}.  Equations
\eqref{eq:verify-we}--\eqref{eq:verify-wi} then give
\[
w_e=D_{\QP}^{(k_e)}\tau+S_{\QP}^{(k_e)}\sigma,
\qquad
w_i=D^{(k_i)}\tau+S^{(k_i)}\sigma.
\]
Substitution into \eqref{eq:theorem-w} is exactly
\eqref{eq:correctedrep}, proving existence.

\emph{4. Uniqueness.}
For a fixed physical field \(u\), the port Cauchy data are fixed.
Formula \eqref{eq:portcoeff} therefore fixes \(A_m^L\) and \(B_m^R\), so the
normalization \eqref{eq:theorem-uh} fixes \(u_h\) uniquely.  Consequently
\((w_e,w_i)\) in \eqref{eq:theorem-w} is fixed.  The uniqueness part of
Lemma~\ref{lem:synthesis}, equivalently \(\Ker\mathcal J=\{0\}\), fixes
\((\tau,\sigma)\) uniquely.  Finally, the modewise Wronskian calculation in
Section~\ref{sec:rayleigh} shows that the Rayleigh coefficients themselves are
unique.  This proves every uniqueness assertion in the theorem.
\end{proof}

\subsection{Why there is no additional complementary resonance here}

Ordinary interior Dirichlet or Neumann eigenvalues can make a pure single- or
double-layer representation nonunique; Kress gives the classical example in
\cite[Section 1]{Kress1991}.  The common-density Muller construction avoids
those defects by coupling the two potentials and using the swapped transmission
problem.  In the present mixed QP/free ansatz, that swapped problem is the
standard full-plane problem \eqref{eq:complement}.  Its homogeneous radiating
solution is zero, so reversing the jump construction cannot produce a nonzero
element of \(\Ker\mathcal J\).  In particular,
\[
\Ker\mathcal J=\{0\}
\]
under the hypotheses of Theorem~\ref{thm:corrected}.  A complementary periodic
resonance would become relevant only for a different formulation using QP
kernels on both physical sides.

\subsection{Physical guided modes are not the same defect}

A physical guided mode has \(u_h=0\) and a nonzero density satisfying
\[
A_{\QP}\eta=0.
\]
This is precisely what the TEP eigenvalue method seeks.  Injectivity of
\(\mathcal J\), established through the free-space complementary problem,
shows that this density generates a nonzero physical field.  Therefore
\(\Ker A_{\QP}\neq\{0\}\) need not imply coordinate redundancy: \(A_{\QP}\)
records only the interface mismatch, whereas \(\mathcal J\) records the whole
field.  Confusing \(\Ker A_{\QP}\) with \(\Ker\mathcal J\) is the central
uniqueness trap.

\section{Wood anomalies and other exceptional cases}

At \(\gamma_{m_0}=0\), three failures occur simultaneously.

\begin{enumerate}
\item The Green series \eqref{eq:qpgreen} contains \(1/\gamma_{m_0}\) and is
undefined without regularization.
\item The two exponentials for order \(m_0\) coalesce, while the ODE requires
the two-dimensional span of \(\psi_{m_0}\) and \(x\psi_{m_0}\).
\item Incoming/outgoing classification of the threshold order is not selected
by the sign of \(\gamma_m\).  A limiting-absorption or generalized-mode rule is
needed.
\end{enumerate}

Thus merely saying ``the proof is ill-conditioned at Wood'' is too weak: the
formula in the conjecture is structurally incomplete there.  A threshold-safe
theorem must change both the Green kernel and the homogeneous basis.  Such a
change is outside this task because it changes the formulation rather than
proving the stated one.

If \(\Im k>0\), the modal branch has strict decay and energy identities usually
remove real-axis guided modes.  This gives a useful regularized theory and a
limiting-absorption route, but the limit back to real \(k\) can be singular at
thresholds or embedded guided modes.

If \(\Om\) touches a vertical port, the homogeneous collar disappears and
\eqref{eq:portcoeff} cannot be used.  If it crosses the quasiperiodic seam, one
must use periodic charts and include the continued interface.  Neither case is
covered by Theorem~\ref{thm:corrected}.

\section{Consequences for the DtN--Muller formulation}

The later draft system uses
\[
A_{\QP}\eta+B_\Sig\xi=0
\]
together with port trace equations.  The present analysis gives five concrete
requirements.

\begin{enumerate}
\item The vector \(\xi=(a,b)\) must be interpreted as left-incoming and
right-incoming homogeneous coordinates, not as an arbitrary duplicate of the
outgoing layer far field.
\item The coefficient domain should be
\(\ell^2_{1/2,\beta}\times\ell^2_{1/2,\beta}\), with the finite-dimensional
truncation viewed only as an approximation.
\item At the continuous level, injectivity of the mixed-kernel synthesis map
ensures that a nonzero density produces a nonzero physical field.  A numerical
implementation should still check for discretization-induced near-nullvectors.
\item Physical singularity of \(A_{\QP}\) is not automatically spurious; it
can be the desired guided mode.  The field synthesis map, not the mismatch row
alone, distinguishes physical and coordinate nullspaces.
\item Near Wood thresholds the standard QP Green matrix and Rayleigh basis
must not be trusted.  Excluding a numerical neighborhood is a practical choice,
but a theorem requires a regularized threshold formulation.
\end{enumerate}

The unconditional direct Green representation in Theorem
\ref{thm:directgreen} remains a useful independent check: it uses physical
Cauchy data as primary unknowns and retains the vertical boundary term
explicitly, at the cost of a different operator system.

\section{Final classification and remaining open work}

\begin{center}
\begin{tabularx}{\textwidth}{>{\raggedright\arraybackslash}p{0.30\textwidth}X}
\toprule
Claim & Status \\
\midrule
Empty-cell non-Wood Rayleigh representation & Proved, including the natural
half-order coefficient norm.\\
Top/bottom Green terms cancel & Proved from opposite normals and reciprocal
quasiperiodic phases.\\
Vertical terms are a homogeneous background & Proved inside the open cell;
they have the Rayleigh representation of Section~\ref{sec:rayleigh}.\\
Direct physical-trace layer representation & Proved in Theorem
\ref{thm:directgreen}; interior/exterior signs are opposite.\\
Draft common-density representation & Proved under the explicit regularity and
non-Wood assumptions of Theorem~\ref{thm:corrected}.\\
Uniqueness of Rayleigh coefficients & Proved away from Wood thresholds after
incoming normalization.\\
Uniqueness of common densities & Proved using uniqueness of the standard
full-plane swapped transmission problem.\\
Original Conjecture 1 as written & Too imprecise to be a theorem; after its
spaces, port normalization, signs, and non-Wood condition are made explicit,
it is replaced by Theorem~\ref{thm:corrected}.\\
\bottomrule
\end{tabularx}
\end{center}

The remaining work concerns genuine exceptional sets and numerical
verification.  At Wood parameters one must construct a threshold-safe QP Green
kernel and generalized Rayleigh basis.  One must also verify that discrete
trace and layer maps preserve the continuous injectivity established above.
The mixed QP/free representation has no separate complementary cylinder
spectrum.

\begin{thebibliography}{99}

\bibitem{BarnettGreengard2010}
A. H. Barnett and L. Greengard,
``A new integral representation for quasiperiodic fields and its application
to two-dimensional band structure calculations,''
\emph{Journal of Computational Physics}, 229(19):6898--6914, 2010.
\href{https://doi.org/10.1016/j.jcp.2010.05.029}{doi:10.1016/j.jcp.2010.05.029};
\href{https://arxiv.org/abs/1001.5464}{arXiv:1001.5464}.

\bibitem{ColtonKress1983}
D. Colton and R. Kress,
\emph{Integral Equation Methods in Scattering Theory},
Classics in Applied Mathematics 72, Society for Industrial and Applied
Mathematics, Philadelphia, 2013; reprint of the original Wiley edition, 1983.
\href{https://doi.org/10.1137/1.9781611973167}{doi:10.1137/1.9781611973167}.

\bibitem{DominguezLyonTurc2016}
V. Dominguez, M. Lyon, and C. Turc,
``Well-posed boundary integral equation formulations and Nystrom
discretizations for the solution of Helmholtz transmission problems in
two-dimensional Lipschitz domains,'' 2016.
\href{https://arxiv.org/abs/1509.04415}{arXiv:1509.04415v2}.

\bibitem{HiptmairMoiolaSpence2022}
R. Hiptmair, A. Moiola, and E. A. Spence,
``Spurious quasi-resonances in boundary integral equations for the Helmholtz
transmission problem,'' \emph{SIAM Journal on Applied Mathematics},
82(4):1446--1469, 2022.
\href{https://doi.org/10.1137/21M1447052}{doi:10.1137/21M1447052};
\href{https://arxiv.org/abs/2109.08530}{arXiv:2109.08530}.

\bibitem{Kress1991}
R. Kress,
``Boundary integral equations in time-harmonic acoustic scattering,''
\emph{Mathematical and Computer Modelling}, 15(3--5):229--243, 1991.
\href{https://doi.org/10.1016/0895-7177(91)90068-I}{doi:10.1016/0895-7177(91)90068-I}.

\bibitem{Linton1998}
C. M. Linton,
``The Green's function for the two-dimensional Helmholtz equation in periodic
domains,'' \emph{Journal of Engineering Mathematics}, 33:377--402, 1998.
\href{https://doi.org/10.1023/A:1004377501747}{doi:10.1023/A:1004377501747}.

\bibitem{McLean2000}
W. McLean,
\emph{Strongly Elliptic Systems and Boundary Integral Equations},
Cambridge University Press, 2000.
ISBN 978-0-521-66332-8.

\end{thebibliography}

\appendix
\section{A modewise uniqueness calculation}

For one non-Wood order, suppose
\[
a\ee^{\ii\gamma(x-X_L)}+b\ee^{-\ii\gamma(x-X_R)}=0
\quad\text{for }X_L<x<X_R.
\]
Evaluating this identity and its derivative at any point gives
\[
\begin{bmatrix}
\ee^{\ii\gamma(x-X_L)}&\ee^{-\ii\gamma(x-X_R)}\\
\ii\gamma\ee^{\ii\gamma(x-X_L)}&
-\ii\gamma\ee^{-\ii\gamma(x-X_R)}
\end{bmatrix}
\begin{bmatrix}a\\b\end{bmatrix}=0.
\]
The determinant is
\(-2\ii\gamma\ee^{\ii\gamma(X_R-X_L)}\neq0\), so \(a=b=0\).
This is coefficient uniqueness.  It says nothing about density uniqueness,
which is controlled by the entirely different map \(\mathcal J\).

\section{Why opposing horizontal Green terms cancel}

Let the source variable be \(z=(z_x,z_y)\).  If the field has phase
\(u(z_x,d)=\alpha u(z_x,0)\), \(\alpha=\ee^{\ii\beta d}\), then the Green
kernel used against it has inverse source phase
\(G(x;z_x,d)=\alpha^{-1}G(x;z_x,0)\).  The same holds for the corresponding
source-normal derivatives.  Since the outward normal is \(+e_y\) at the top
and \(-e_y\) at the bottom, substituting the two phase relations into
\[
\int_{\Gamma_{\rm top}\cup\Gamma_{\rm bot}}
  (G\partial_n u-u\partial_nG)\,ds
\]
shows term-by-term cancellation.  This is why no extra horizontal-boundary
unknown is required when the kernel and field use matching quasiperiodicity.

\section{Operator-level uniqueness dictionary}

\begin{align*}
\text{PDE field uniqueness}
&\neq \text{representation uniqueness},\\
\Ker A_{\QP}
&=\text{densities producing zero interface mismatch},\\
\Ker\mathcal J
&=\text{densities producing zero fields on both physical sides},\\
\Ker A_{\QP}\setminus\Ker\mathcal J
&=\text{physical guided-mode densities (under the stated hypotheses)},\\
\Ker\mathcal J
&=\{0\}\quad\text{by full-plane complementary uniqueness},\\
\gamma_m=0
&\Rightarrow\text{failure of the displayed Rayleigh coordinate basis}.
\end{align*}

This dictionary is the shortest way to avoid conflating the physical mismatch
nullspace, the synthesis map, and the separate Wood-threshold failure.

\end{document}
